\documentclass[preview]{standalone}

\usepackage{amsmath}
\usepackage{amssymb}
\usepackage{stellar}
\usepackage{definitions}
\usepackage{bettelini}

\begin{document}

\title{Stellar}
\id{linearalgebra-characteristic-polynomial}
\genpage

\section{Trace}

\begin{snippetdefinition}{matrix-trace-definition}{Trace of a matrix}
    Let \(\mathbb{F}\) be a \field, and let
    \(A=(a_{ij})\in\matrices_{n\times n}(\mathbb{F})\). The
    \emph{trace} of \(A\) is
    \[
        \operatorname{tr}(A)=\sum_{i=1}^{n}a_{ii}.
    \]
\end{snippetdefinition}

\begin{snippetproposition}{trace-cyclic-property}{Cyclic property of the trace}
    Let \(\mathbb{F}\) be a \field. If
    \(A,B\in\matrices_{n\times n}(\mathbb{F})\), then
    \[
        \operatorname{tr}(AB)=\operatorname{tr}(BA).
    \]
\end{snippetproposition}

\begin{snippetproof}{trace-cyclic-property-proof}{trace-cyclic-property}{Cyclic property of the trace}
    If \(A=(a_{ij})\) and \(B=(b_{ij})\), then
    \[
        \operatorname{tr}(AB)
        =
        \sum_{i=1}^{n}\sum_{j=1}^{n}a_{ij}b_{ji}.
    \]
    Reindexing the same finite sum gives
    \[
        \sum_{i=1}^{n}\sum_{j=1}^{n}a_{ij}b_{ji}
        =
        \sum_{j=1}^{n}\sum_{i=1}^{n}b_{ji}a_{ij}
        =
        \operatorname{tr}(BA).
    \]
\end{snippetproof}

\begin{snippetcorollary}{similar-matrices-have-same-trace}{Similar matrices have the same trace}
    Let \(\mathbb{F}\) be a \field. If
    \(B,C\in\matrices_{n\times n}(\mathbb{F})\) are similar, then
    \[
        \operatorname{tr}(B)=\operatorname{tr}(C).
    \]
\end{snippetcorollary}

\begin{snippetproof}{similar-matrices-have-same-trace-proof}{similar-matrices-have-same-trace}{Similar matrices have the same trace}
    Since \(B\) and \(C\) are similar, there exists an invertible
    \(A\in\operatorname{GL}(n,\mathbb{F})\) such that
    \[
        B=A^{-1}CA.
    \]
    By the cyclic property,
    \[
        \operatorname{tr}(B)
        =
        \operatorname{tr}(A^{-1}CA)
        =
        \operatorname{tr}(CAA^{-1})
        =
        \operatorname{tr}(C).
    \]
\end{snippetproof}

\begin{snippetdefinition}{endomorphism-trace-definition}{Trace of an endomorphism}
    Let \(V\) be a finite-dimensional \vectorspace over a \field
    \(\mathbb{F}\), and let \(f\in\operatorname{End}(V)\). The
    \emph{trace of \(f\)} is
    \[
        \operatorname{tr}(f)=\operatorname{tr}\bigl(M(f,\mathcal{B},\mathcal{B})\bigr),
    \]
    where \(\mathcal{B}\) is any \basis of \(V\).
\end{snippetdefinition}

\begin{snippetproof}{endomorphism-trace-definition-proof}{endomorphism-trace-definition}{Trace of an endomorphism}
    If \(\mathcal{B}\) and \(\mathcal{E}\) are two \basis[bases] of \(V\),
    then the associated \matrix[matrices] \(M(f,\mathcal{B},\mathcal{B})\) and
    \(M(f,\mathcal{E},\mathcal{E})\) are similar. \similarmatrices[Similar matrices] have the
    same trace, so the definition does not depend on the choice of
    \(\mathcal{B}\).
\end{snippetproof}

\section{Characteristic polynomial}

\begin{snippetdefinition}{matrix-characteristic-polynomial-definition}{Characteristic polynomial of a matrix}
    Let \(\mathbb{F}\) be a \field, and let
    \(A\in\matrices_{n\times n}(\mathbb{F})\). The
    \emph{characteristic polynomial} of \(A\) is
    \[
        p_A(t)=\det(t\identmatrix{n}-A)\in\mathbb{F}[t].
    \]
\end{snippetdefinition}

\begin{snippetdefinition}{characteristic-polynomial-definition}{Characteristic polynomial of an endomorphism}
    Let \(V\) be a finite-dimensional \vectorspace over a \field
    \(\mathbb{F}\), with \(\lineardim V=n\), and let
    \(f\in\operatorname{End}(V)\). The \emph{characteristic polynomial} of
    \(f\) is
    \[
        p_f(t)=\det(t\identmatrix{n}-M(f,\mathcal{B},\mathcal{B})),
    \]
    where \(\mathcal{B}\) is any \basis of \(V\).
\end{snippetdefinition}

\begin{snippetproof}{characteristic-polynomial-definition-proof}{characteristic-polynomial-definition}{Characteristic polynomial of an endomorphism}
    If \(\mathcal{B}\) and \(\mathcal{E}\) are two \basis[bases] of \(V\),
    then
    \[
        M(f,\mathcal{E},\mathcal{E})
        =
        P^{-1}M(f,\mathcal{B},\mathcal{B})P
    \]
    for some invertible \matrix \(P\). Hence
    \[
        t\identmatrix{n}-M(f,\mathcal{E},\mathcal{E})
        =
        P^{-1}
        \bigl(t\identmatrix{n}-M(f,\mathcal{B},\mathcal{B})\bigr)
        P.
    \]
    \similarmatrices[Similar matrices] have the same determinant, so the \polynomial is
    independent of the chosen \basis.
\end{snippetproof}

\begin{snippetproposition}{characteristic-polynomial-properties}{Basic properties of the characteristic polynomial}
    Let \(V\) be a finite-dimensional \vectorspace over a \field
    \(\mathbb{F}\), with \(\lineardim V=n\), and let
    \(f\in\operatorname{End}(V)\). Then:
    \begin{enumerate}
        \item \(p_f(t)\) is a monic \polynomial of degree \(n\);
        \item the coefficient of \(t^{n-1}\) is \(-\operatorname{tr}(f)\);
        \item \(p_f(0)=(-1)^n\det(f)\);
        \item \(\lambda\in\mathbb{F}\) is an \eigenvalue of \(f\) if and only if \(p_f(\lambda)=0\).
    \end{enumerate}
\end{snippetproposition}

\begin{snippetproof}{characteristic-polynomial-properties-proof}{characteristic-polynomial-properties}{Basic properties of the characteristic polynomial}
    Choose a \basis \(\mathcal{B}\) of \(V\), and put
    \[
        A=M(f,\mathcal{B},\mathcal{B})=(a_{ij}).
    \]
    Then
    \[
        p_f(t)=\det(t\identmatrix{n}-A).
    \]
    In the matrix \(t\identmatrix{n}-A\), the diagonal entries are
    \(t-a_{11},\ldots,t-a_{nn}\). The term of degree \(n\) is obtained by
    choosing \(t\) from every diagonal entry, so \(p_f(t)\) is monic of
    degree \(n\).

    The coefficient of \(t^{n-1}\) is obtained by choosing one constant
    diagonal term and \(t\) from all remaining diagonal terms. Hence it is
    \[
        -(a_{11}+\cdots+a_{nn})
        =
        -\operatorname{tr}(A)
        =
        -\operatorname{tr}(f).
    \]
    Moreover,
    \[
        p_f(0)=\det(-A)=(-1)^n\det(A)=(-1)^n\det(f).
    \]

    Finally,
    \[
        p_f(\lambda)=0
        \Longleftrightarrow
        \det(\lambda\identmatrix{n}-A)=0.
    \]
    This is equivalent to the existence of a nonzero vector in
    \[
        \ltkernel(\lambda\identityfunc_V-f),
    \]
    which is equivalent to \(\lambda\) being an \eigenvalue of \(f\).
\end{snippetproof}

\section{Polynomial roots}

\begin{snippettheorem}{polynomial-division-over-field-theorem}{Polynomial division over a field}
    Let \(\mathbb{F}\) be a \field, and let
    \(u,v\in\mathbb{F}[t]\), with \(v\neq0\). There exist unique
    \polynomial[polynomials] \(q,r\in\mathbb{F}[t]\) such that
    \[
        u=vq+r
    \]
    and either \(r=0\) or
    \[
        \polynomialdeg r<\polynomialdeg v.
    \]
\end{snippettheorem}

\begin{snippetproof}{polynomial-division-over-field-theorem-proof}{polynomial-division-over-field-theorem}{Polynomial division over a field}
    For uniqueness, suppose
    \[
        u=vq+r=vq'+r',
    \]
    with \(r,r'\) either zero or of degree smaller than \(\polynomialdeg v\).
    Then
    \[
        v(q-q')=r'-r.
    \]
    If \(q\neq q'\), the left-hand side has degree at least
    \(\polynomialdeg v\), while the right-hand side has degree smaller than
    \(\polynomialdeg v\), a contradiction. Hence \(q=q'\), and then
    \(r=r'\).

    For existence, if \(\polynomialdeg u<\polynomialdeg v\), take
    \(q=0\) and \(r=u\). Otherwise, subtract from \(u\) a suitable multiple
    of \(v\) that cancels the highest-degree term of \(u\). Repeating this
    process strictly lowers the degree at each step, so after finitely many
    steps one obtains a remainder \(r\) of degree smaller than
    \(\polynomialdeg v\), or \(r=0\).
\end{snippetproof}

\begin{snippetcorollary}{polynomial-root-linear-factor-corollary}{Roots and linear factors}
    Let \(\mathbb{F}\) be a \field, let \(p\in\mathbb{F}[t]\), and let
    \(\lambda\in\mathbb{F}\). Then
    \[
        p(\lambda)=0
        \iff
        (t-\lambda)\text{ divides }p(t).
    \]
\end{snippetcorollary}

\begin{snippetproof}{polynomial-root-linear-factor-corollary-proof}{polynomial-root-linear-factor-corollary}{Roots and linear factors}
    Divide \(p\) by \(t-\lambda\). There exist
    \(q,r\in\mathbb{F}[t]\) such that
    \[
        p(t)=(t-\lambda)q(t)+r(t),
    \]
    where \(r\) is constant. Evaluating at \(t=\lambda\) gives
    \[
        p(\lambda)=r.
    \]
    Hence \(p(\lambda)=0\) if and only if \(r=0\), which is equivalent to
    divisibility by \(t-\lambda\).
\end{snippetproof}

\begin{snippetcorollary}{polynomial-distinct-roots-bound}{Number of roots of a polynomial}
    Let \(\mathbb{F}\) be a \field. A nonzero \polynomial
    \(p\in\mathbb{F}[t]\) of degree \(n\) has at most \(n\) distinct roots
    in \(\mathbb{F}\).
\end{snippetcorollary}

\begin{snippetproof}{polynomial-distinct-roots-bound-proof}{polynomial-distinct-roots-bound}{Number of roots of a polynomial}
    Each distinct root \(\lambda\) gives a distinct linear factor
    \(t-\lambda\). Removing one such factor lowers the degree by \(1\).
    Thus a nonzero \polynomial of degree \(n\) cannot have more than \(n\)
    distinct roots.
\end{snippetproof}

\includesnpt{fundamental-theorem-of-algebra}

\begin{snippetcorollary}{complex-polynomial-factorization-corollary}{Factorization over \(\mathbb{C}\)}
    Every nonzero \polynomial \(p\in\complexnumbers[t]\) of degree \(n\) can
    be written as
    \[
        p(t)=c(t-\lambda_1)\cdots(t-\lambda_n),
    \]
    where \(c\in\complexnumbers\setminus\{0\}\) and
    \(\lambda_1,\ldots,\lambda_n\in\complexnumbers\), with roots repeated
    according to multiplicity.
\end{snippetcorollary}

\begin{snippetproof}{complex-polynomial-factorization-corollary-proof}{complex-polynomial-factorization-corollary}{Factorization over \(\mathbb{C}\)}
    This follows from the fundamental theorem of algebra by repeatedly
    factoring out a linear factor corresponding to a root.
\end{snippetproof}

\begin{snippetproposition}{real-polynomial-factorization-proposition}{Factorization over \(\mathbb{R}\)}
    Every nonzero \polynomial \(p\in\realnumbers[t]\) can be written as
    \[
        p(t)=
        a(t-\lambda_1)\cdots(t-\lambda_h)q_1(t)\cdots q_r(t),
    \]
    where \(a\in\realnumbers\setminus\{0\}\),
    \(\lambda_1,\ldots,\lambda_h\in\realnumbers\), and every \(q_i\) is a
    monic quadratic \polynomial in \(\realnumbers[t]\) with no real roots.
\end{snippetproposition}

\begin{snippetproof}{real-polynomial-factorization-proposition-proof}{real-polynomial-factorization-proposition}{Factorization over \(\mathbb{R}\)}
    Regard \(p\) as a complex polynomial. By the complex factorization
    theorem, it splits into linear factors over \(\complexnumbers\). Since
    the coefficients of \(p\) are real, every non-real root \(z\) occurs
    together with its complex conjugate \(\overline{z}\). Pairing conjugate
    roots gives factors
    \[
        (t-z)(t-\overline{z})
        =
        t^2-2\operatorname{Re}(z)t+|z|^2,
    \]
    which have real coefficients and no real roots. The remaining roots are
    real and give the linear factors.
\end{snippetproof}

\begin{snippetexample}{real-polynomial-factorization-example}{A real factorization}
    In \(\realnumbers[t]\),
    \[
        t^4+1
        =
        (t^2+\sqrt{2}t+1)(t^2-\sqrt{2}t+1).
    \]
    Both quadratic factors have no real roots.
\end{snippetexample}

\begin{snippetexample}{rational-polynomials-without-roots-example}{Polynomials without rational roots}
    For every \(k\geq1\), the \polynomial
    \[
        t^{2k}-2
    \]
    has no root in \(\rationalnumbers\). Indeed, a rational root would be a
    rational number \(q\) such that
    \[
        q^{2k}=2,
    \]
    which is impossible. Indeed, if \(q=a/b\) with \(a,b\in\integers\),
    \(b\neq0\), and \(\gcd(a,b)=1\), then
    \[
        a^{2k}=2b^{2k}.
    \]
    Thus \(a\) is even. Writing \(a=2c\), one gets
    \[
        2^{2k}c^{2k}=2b^{2k},
    \]
    so \(b^{2k}\) is even, and hence \(b\) is even. This contradicts
    \(\gcd(a,b)=1\).
\end{snippetexample}

\end{document}
